% ================================================================
%  CHAPTER 2: LITERATURE REVIEW
% ================================================================
\chapter{Literature Review}
\label{chap:lit_review}

This chapter provides a comprehensive review of the scientific literature
underpinning this thesis. Section~\ref{sec:egfr_bio} reviews the molecular
biology of EGFR and its role in NSCLC oncogenesis. Section~\ref{sec:tki_gen}
discusses the three generations of EGFR TKIs and their clinical limitations.
Section~\ref{sec:chembl_db} describes the ChEMBL bioactivity database and
associated data quality challenges. Section~\ref{sec:molrep} reviews molecular
representation methods for machine learning. Section~\ref{sec:gnn_arch} surveys the GNN architectures employed in this study.
Section~\ref{sec:nepali_plants} reviews the Nepali medicinal plants and natural
product databases used for virtual screening.
Section~\ref{sec:vs_lit} reviews virtual screening methodologies.
Section~\ref{sec:research_gap} synthesises the identified research gaps
that motivate this thesis.

% ----------------------------------------------------------------
\section{EGFR Kinase Biology and NSCLC}
\label{sec:egfr_bio}
% ----------------------------------------------------------------

\subsection{Molecular Architecture of EGFR}

The Epidermal Growth Factor Receptor (EGFR), also designated ErbB1 or HER1,
is a 134 kDa single-pass transmembrane receptor tyrosine kinase (RTK) encoded
by the \textit{EGFR} gene on chromosome 7p11.2. EGFR is the founding member of
the ErbB receptor family (ErbB1--ErbB4), which collectively orchestrates
fundamental cellular processes including proliferation, differentiation, survival,
and motility \parencite{hynes2005erbb}. Structurally, EGFR comprises a glycosylated \textbf{extracellular domain} (I--IV)
that binds EGF-family ligands and mediates dimerisation, a single-pass
\textbf{transmembrane} helix, a regulatory \textbf{juxtamembrane} segment, and a
bilobed intracellular \textbf{kinase domain}. The ATP-binding cleft---formed at the
N-/C-lobe interface and lined by the hinge, P-loop, activation loop, and DFG
motif---is the target of all small-molecule EGFR inhibitors, and a carboxy-terminal
tail bears the autophosphorylation sites (Tyr1068, Tyr1173, and others) that recruit
downstream effectors.

Under normal conditions EGFR is autoinhibited; ligand binding drives homo- or
heterodimerisation with other ErbB members (HER2--HER4), relieving autoinhibition and
triggering \textit{trans}-autophosphorylation that activates the RAS--RAF--MEK--ERK
(MAPK), PI3K--AKT--mTOR, STAT3/5, and PLC-$\gamma$ signalling cascades---collectively
driving proliferation, survival, and differentiation \parencite{siegel2023cancer}.

\subsection{Oncogenic EGFR Mutations in NSCLC}

In non-small cell lung cancer (NSCLC), somatic gain-of-function mutations within
the EGFR kinase domain render the receptor constitutively active in a
ligand-independent manner, producing continuous, dysregulated stimulation of the
pro-survival and proliferative pathways described above \parencite{rotow2023egfr}.
This oncogenic signalling drives uncontrolled tumour growth, resistance to
apoptosis, angiogenesis, and metastatic dissemination.

The two most clinically prevalent sensitising EGFR mutations, which together
account for 85--90\% of all EGFR-mutant NSCLC, are:

\begin{itemize}
    \item \textbf{Exon 19 deletions (Ex19del)}: A heterogeneous group of in-frame
    deletions of 9--24 base pairs within exon 19, most commonly removing the
    LREA motif (residues 747--750) from the $\alpha$C-helix of the kinase N-lobe.
    These deletions stabilise the active conformation of the kinase domain by
    repositioning the $\alpha$C-helix, reducing the activation energy for catalysis.
    Ex19del accounts for approximately 45--50\% of EGFR-mutant NSCLC.
    \item \textbf{L858R point mutation}: A leucine-to-arginine substitution at
    position 858 within the activation loop (A-loop) of the kinase domain.
    This mutation disrupts the autoinhibitory network involving Leu858, Asp855,
    and Phe856 of the DFG motif, shifting the equilibrium from the inactive to
    the active kinase conformation. L858R accounts for approximately 40\% of
    EGFR-mutant NSCLC.
\end{itemize}

Rare sensitising mutations ($<$5\% combined) include exon 18 G719X, exon 20
S768I, and exon 21 L861Q. These display differential sensitivity to approved
TKIs and often require alternative treatment strategies.

The prevalence of EGFR mutations in NSCLC exhibits a profound ethnic disparity.
In Caucasian populations, EGFR mutations are found in approximately 10--15\% of
NSCLC cases, whereas in East Asian populations (Chinese, Japanese, Korean, and
related groups), the prevalence reaches 40--50\% \parencite{leonetti2019resistance}.
This epidemiological pattern reflects underlying differences in the frequency of
germline polymorphisms at the \textit{EGFR} locus that predispose to somatic
mutation acquisition. Importantly, this disparity has direct clinical relevance
for South and East Asia, including Nepal, where EGFR-targeted therapy represents
the most impactful precision oncology intervention available, yet access to
diagnostic testing and novel therapeutic agents remains significantly constrained.

% ----------------------------------------------------------------
\section{Three Generations of EGFR TKIs and Resistance Mechanisms}
\label{sec:tki_gen}
% ----------------------------------------------------------------

\subsection{First-Generation Reversible TKIs}

The clinical development of EGFR tyrosine kinase inhibitors began with the
regulatory approval of gefitinib (Iressa, AstraZeneca) in 2003 and erlotinib
(Tarceva, Genentech/Roche) in 2004. Both agents are reversible, ATP-competitive
small molecule inhibitors that occupy the EGFR ATP-binding cleft through a
combination of a hinge-region hydrogen bond (to the backbone amide nitrogen of
Met793), hydrophobic contacts with the ribose pocket, and van der Waals interactions
with the hydrophobic regions I and II of the binding site. Crystal structures of
EGFR bound to gefitinib (PDB: 2ITY) and erlotinib (PDB: 1M17) have confirmed
these binding modes at the atomic level \parencite{stamos2002egfr}.

Landmark clinical trials established the superiority of first-generation TKIs
over platinum-doublet chemotherapy in patients with sensitising EGFR mutations.
The IPASS trial (gefitinib vs. carboplatin/paclitaxel) and EURTAC trial
(erlotinib vs. platinum doublet) demonstrated objective response rates of
55--75\% and median progression-free survival (PFS) of 9--13 months---unprecedented
outcomes for a molecularly unselected metastatic NSCLC population.

However, acquired resistance develops in virtually all patients within 9--13 months.
The most common resistance mechanism (50--60\% of cases) is the \textbf{T790M
gatekeeper mutation} (Thr790Met) in exon 20 of the EGFR kinase domain
\parencite{leonetti2019resistance}. This methionine substitution acts through
two complementary mechanisms: (i) steric clash with the quinazoline scaffold
of first-generation inhibitors due to the bulkier methionine side chain, and
(ii) a 100-fold increase in the binding affinity of the kinase for ATP, effectively
out-competing the reversible TKI. Additional resistance mechanisms include MET
amplification ($\sim$5--15\%), HER2 amplification ($\sim$5--10\%), PIK3CA
mutation ($\sim$5\%), and histological transformation to small cell carcinoma
($\sim$5\%) \parencite{leonetti2019resistance}.

\subsection{Second-Generation Irreversible TKIs}

The T790M-mediated resistance to first-generation TKIs motivated the development
of second-generation irreversible inhibitors: afatinib (Gilotrif, Boehringer
Ingelheim) and dacomitinib (Vizimpro, Pfizer). These agents form a covalent
Michael addition bond with the thiol group of Cys797 at the rim of the ATP-binding
cleft, providing kinetically stable, irreversible target engagement. As pan-ErbB
inhibitors, they also inhibit HER2 and HER4 kinase activity.

Afatinib demonstrated superiority over platinum-doublet chemotherapy in the LUX-Lung
3 and LUX-Lung 6 trials, and superior PFS versus gefitinib in LUX-Lung 7, establishing
it as a first-line option for EGFR-mutant NSCLC. Dacomitinib similarly improved
PFS over gefitinib in the ARCHER 1050 trial. However, the irreversible
pan-ErbB inhibition profile of second-generation agents produces significant
dose-limiting on-target toxicity---primarily grade 3/4 diarrhoea, skin rash,
stomatitis, and paronychia---arising from inhibition of wild-type EGFR in the
gastrointestinal epithelium and skin \parencite{leonetti2019resistance}.
This toxicity profile restricts clinical dosing to sub-therapeutic levels that
are insufficient to overcome T790M-driven resistance in the majority of patients.

\subsection{Third-Generation Covalent TKIs and the C797S Resistance Problem}

The clinical unmet need created by T790M resistance drove the development of
third-generation mutant-selective covalent TKIs, led by osimertinib (Tagrisso,
AstraZeneca; AZD9291). Osimertinib was rationally designed through structure-based
drug design to selectively engage T790M-mutant EGFR while sparing wild-type EGFR.
Its indole-based acrylamide scaffold binds the EGFR ATP cleft and forms an
irreversible covalent bond with Cys797 through a Michael addition, while its
selectivity for mutant over wild-type EGFR is achieved through differential
binding geometries and kinetics rather than steric exclusion
\parencite{rotow2023egfr}.

The AURA3 trial established osimertinib as the standard second-line therapy for
T790M-positive NSCLC after first-generation TKI failure. Subsequently, the
FLAURA trial demonstrated osimertinib's superiority over first-generation TKIs
as first-line therapy (median PFS 18.9 vs. 10.2 months; overall survival 38.6
vs. 31.8 months), establishing it as the unambiguous standard of care for
EGFR-mutant NSCLC regardless of prior treatment. Additional advantages include
excellent CNS penetration ($\sim$75\% brain metastasis response rate) and an
improved safety profile compared to afatinib.

Despite these remarkable clinical achievements, acquired resistance to osimertinib
emerges in virtually all patients within 9--18 months \parencite{leonetti2019resistance}.
The dominant resistance mechanism is the \textbf{C797S mutation} (Cys797Ser),
which substitutes the cysteine residue required for covalent TKI binding with a
serine, abolishing the Michael addition reaction. C797S is detected in 20--40\%
of post-osimertinib biopsies and is of particular clinical importance when it
occurs in \textit{cis} with T790M on the same allele, generating the
\textbf{triple-mutant EGFR} (Ex19del/T790M/C797S and L858R/T790M/C797S) that is
resistant to \textit{all three generations} of approved EGFR TKIs:
\begin{itemize}
    \item First-generation TKIs cannot overcome T790M.
    \item Second-generation TKIs cannot overcome T790M steric clash at therapeutic doses.
    \item Third-generation TKIs cannot bind due to loss of the Cys797 covalent anchor.
\end{itemize}

Additional post-osimertinib resistance mechanisms include acquired MET amplification
($\sim$15\%), HER2 amplification ($\sim$5\%), RAS/RAF/MAP2K mutations ($\sim$10\%),
EGFR amplification, EGFR C797 allosteric mutations (e.g., G796S, L792 mutations),
and histological transformation to SCLC or squamous cell carcinoma
\parencite{rotow2023egfr}. As of 2026, no approved therapeutic agent effectively
addresses triple-mutant EGFR. Investigational fourth-generation strategies---including
reversible, mutant-selective non-covalent inhibitors (e.g., BLU-945) and allosteric
inhibitors that bind outside the ATP pocket (e.g., EAI045)---are in preclinical and
early clinical development, but none is yet approved and the durability of response
remains uncertain. This makes the discovery of structurally novel, non-covalent EGFR
inhibitor chemotypes---including those derived from natural products---the highest
clinical priority in NSCLC drug development.

% ----------------------------------------------------------------
\section{Bioactivity Databases and the ChEMBL Database}
\label{sec:chembl_db}
% ----------------------------------------------------------------

\subsection{Overview of the ChEMBL Database}

The ChEMBL database, maintained by EMBL-EBI, is the largest curated public repository
of bioactive small molecules and their quantitative activities
\parencite{mendez2019chembl}, comprising over 2.3 million compounds and 20 million
bioactivity measurements across more than 14,000 targets, curated primarily from the
medicinal-chemistry literature. The EGFR entry (target ID: CHEMBL203) is among its
most extensively annotated targets, reflecting decades of drug-discovery activity
around this oncology target.

\subsection{Bioactivity Endpoints and \texorpdfstring{pIC$_{50}$}{pIC50} Conversion}

The primary bioactivity endpoint used in this study is the \textbf{half-maximal
inhibitory concentration} (IC$_{50}$), defined as the concentration of a compound
required to inhibit the target activity by 50\% under specified assay conditions.
IC$_{50}$ values are expressed in molar concentration units (typically nanomolar,
nM) and span many orders of magnitude across structurally diverse compounds.

For machine learning purposes, IC$_{50}$ values are routinely converted to the
logarithmic \textbf{pIC$_{50}$} scale by taking the negative base-ten logarithm of
the molar concentration, so that higher pIC$_{50}$ values correspond to more potent
compounds. This transformation compresses the large dynamic range of IC$_{50}$
values (spanning 6--8 orders of magnitude) into a more tractable numerical range
(approximately 4--11 for drug-like compounds), improves the normality of the
activity distribution, and is the standard practice in QSAR modelling.
A threshold of pIC$_{50} \geq 6.0$ (corresponding to an IC$_{50}$ of $\leq 1\,\mu$M)
is the conventional medicinal chemistry criterion for classifying a compound as a
``hit'' or ``active'' in kinase inhibitor discovery \parencite{wu2018moleculenet}.

\subsection{Data Quality Challenges}

Utilising ChEMBL data for machine learning presents several important quality
challenges that must be addressed during data curation \parencite{wu2018moleculenet}:

\begin{itemize}
    \item \textbf{Assay heterogeneity}: IC$_{50}$ values from hundreds of laboratories
    use different assay formats and conditions, which can shift absolute values by
    1--2 orders of magnitude and add label noise.
    \item \textbf{Censored measurements}: Many values are censored ($>$ or $<$);
    standard practice retains only exact (``='') measurements.
    \item \textbf{Replicate measurements}: Repeated tests of the same compound are
    aggregated by the median pIC$_{50}$ to limit outlier influence.
    \item \textbf{Structural redundancy}: Salt forms and notation differences are
    resolved by RDKit canonicalisation and largest-fragment salt stripping.
    \item \textbf{Activity cliffs}: Structurally similar compounds with very different
    activities challenge both fingerprint- and graph-based models.
\end{itemize}

% ----------------------------------------------------------------
\section{Molecular Representation for Machine Learning}
\label{sec:molrep}
% ----------------------------------------------------------------

\subsection{Hand-Crafted Fingerprint Descriptors}

Molecular fingerprints are fixed-length binary or integer vectors that encode
the presence or count of predefined structural features within a molecule.
They represent the dominant molecular representation paradigm in traditional
cheminformatics and remain strong baselines for bioactivity prediction tasks.

\subsubsection{Extended Connectivity Fingerprints (ECFP)}

The Extended Connectivity Fingerprint (ECFP) \parencite{rogers2010ecfp}, also
known as the Morgan fingerprint (after its algorithmic basis in the Morgan
algorithm for canonical atom numbering), is the most widely used molecular
fingerprint in cheminformatics. ECFP encodes circular atomic neighbourhoods
through an iterative procedure:

\begin{enumerate}
    \item Each atom is assigned an initial integer identifier based on six atomic
    invariants (atomic number, degree, number of hydrogens, charge, isotope,
    and ring membership).
    \item At each iteration $r$ (radius), each atom's identifier is updated by
    hashing together its current identifier with the identifiers of its immediate
    neighbours (corresponding to atoms within $r$ bonds).
    \item All identifiers generated at each iteration up to radius $r$ are
    collected and folded into a bit vector of specified length (typically 1024
    or 2048 bits) using a hash function.
\end{enumerate}

ECFP4 (radius = 2, 2048 bits) is the de facto standard, encoding circular
substructures up to 4 bonds in diameter. Its advantages include computational
efficiency ($O(N)$ in molecular size), interpretability of individual bits as
specific substructural features, and strong empirical performance across diverse
QSAR benchmarks. Limitations include: (i) fixed-length encoding discards
variable-size topological patterns not captured within the predefined radius;
(ii) bit collisions in folded fingerprints reduce information content;
(iii) bond-level information (bond type, conjugation, stereochemistry) is only
coarsely encoded; and (iv) the representation is not differentiable, precluding
end-to-end learning of task-optimal molecular representations.

\subsubsection{Other Descriptor Types}

Beyond fingerprints, molecular property prediction has traditionally relied on
physicochemical descriptors (e.g., RDKit 2D and 3D descriptors, MACCS keys,
DRAGON descriptors), shape-based 3D pharmacophore fingerprints, and continuous
molecular embeddings derived from SMILES string language models.

For the purposes of this study, ECFP4 fingerprints are used for both traditional
machine learning baselines. \textbf{Random Forest (RF)} is an ensemble of 500
decision trees with balanced class weighting and \texttt{max\_features='sqrt'},
providing a well-calibrated, sample-efficient baseline for medium-sized bioactivity
datasets. The \textbf{Multi-Layer Perceptron (MLP)} baseline is a three-layer
neural network (2048$\to$512$\to$128$\to$1) operating on the same ECFP4 inputs,
trained with BCEWithLogitsLoss and early stopping. The MLP bridges the gap between
fingerprint-based and graph-based deep learning, allowing direct attribution of
performance differences to the graph representation rather than to deep learning
capacity in general.

\subsection{Graph-Based Molecular Representations}

The fundamental insight motivating the graph-based approach is that a molecule
is inherently a labelled graph: atoms are nodes with associated chemical properties,
and bonds are edges with associated chemical characteristics. This isomorphism
allows molecules to be fed directly into graph processing algorithms without
information-losing pre-encoding steps \parencite{gilmer2017neural}.

A molecular graph is defined by three components:
\begin{itemize}
    \item A set of \textbf{nodes}, where each node represents one heavy atom and
    carries a node feature vector describing that atom's chemical properties.
    \item A set of \textbf{edges}, where each edge represents one chemical bond and
    carries an edge feature vector describing that bond's properties.
    \item The \textbf{bonding connectivity}, recording which atoms are bonded to
    which.
\end{itemize}

The dimensionality and content of the node and edge feature vectors directly
determine the chemical information available to the model. In this work, a
comprehensive 29-dimensional node feature vector and 12-dimensional edge feature
vector were designed to capture the chemical information most relevant to kinase
inhibitor classification, as detailed in Chapter 3.

\subsection{Graph Neural Networks: The Message-Passing Framework}

The general message-passing neural network (MPNN) framework \parencite{gilmer2017neural}
provides a unified description of most GNN architectures. In this framework, each
atom maintains a hidden representation that is refined layer by layer through two
operations, illustrated in Figure~\ref{fig:msgpass}:

\textbf{Aggregation}: each atom collects ``messages'' from its bonded neighbours,
where each message is a learned function of the central atom's representation, the
neighbour's representation, and the connecting bond features. The messages are then
combined using a permutation-invariant operator (sum, mean, or max) so that the
result does not depend on the arbitrary ordering of neighbours.

\textbf{Update}: each atom updates its own hidden representation by combining its
previous state with the aggregated message, typically through a small multilayer
perceptron or a gated recurrent unit.

After several message-passing layers, a \textbf{readout function} aggregates all
the atom-level representations into a single graph-level embedding, which is then
passed to a task-specific head (for example, a fully connected classifier for
binary activity prediction).

\begin{figure}[htbp]
\centering
\includegraphics[width=0.62\textwidth]{gnn_message_passing.pdf}
\caption{One step of GNN message passing: atom $v$ aggregates learned messages from its bonded neighbours $u_1,u_2,u_3$.}
\label{fig:msgpass}
\end{figure}

The key advantages of the MPNN framework for molecular property prediction are:
(i) the representation is learned end-to-end from bioactivity data, requiring
no expert-designed feature engineering; (ii) arbitrary atom- and bond-level
information can be incorporated via the node and edge feature vectors; (iii) the
framework naturally handles molecules of varying size without fixed-length encoding;
and (iv) the learned representations capture both local chemical environments
(from lower layers) and global molecular topology (from higher layers through
the readout).

% ----------------------------------------------------------------
\section{Graph Neural Network Architectures for Drug Discovery}
\label{sec:gnn_arch}
% ----------------------------------------------------------------

\subsection{Graph Convolutional Network (GCN)}

The Graph Convolutional Network (GCN) \parencite{kipf2017gcn} introduced the
concept of spectral graph convolution through a computationally tractable
first-order approximation. At each layer, every atom's representation is updated by
taking a degree-normalised average of the (linearly transformed) representations of
its bonded neighbours together with its own, followed by a non-linear activation.
The normalisation uses the atom degrees (with self-loops added) so that atoms with
many bonds do not dominate the aggregation.

The GCN architecture has several notable limitations in the context of molecular
property prediction. First, the degree-normalised aggregation weights all
neighbouring atoms equally, without learning which chemical neighbours contribute
most to a given property. Second, standard GCN implementations use node features
only and do not incorporate edge features, discarding bond-level chemical
information. Third, GCN is limited in its expressive power: it cannot distinguish
certain pairs of non-isomorphic graph structures that differ only in their
higher-order topology, limiting its ability to discriminate complex pharmacophoric
patterns \parencite{xu2019gin}. Despite these limitations, GCN provides a strong
and computationally efficient baseline and demonstrated competitive AUROC
performance in the TDC ADMET benchmarks \parencite{huang2021therapeutics}.

\subsection{Graph Attention Network (GAT)}

The Graph Attention Network (GAT) \parencite{velickovic2018gat} replaces the GCN's
fixed degree-based averaging with learnable attention coefficients: each neighbour's
contribution to an atom's update is weighted by a learned score computed from the two
atoms' representations. Several parallel attention ``heads'' are used---concatenated
in intermediate layers and averaged in the final one---and the GATv2 extension
\parencite{brody2022attentive} fixes the ``static attention'' limitation of the
original formulation by conditioning the weights dynamically on both the query and key
nodes. This lets the network focus on the atoms most relevant to activity, such as the
hinge-binding nitrogens of the quinazoline pharmacophore over solvent-exposed
substituents.

\subsection{Graph Isomorphism Network (GIN)}

The Graph Isomorphism Network (GIN) \parencite{xu2019gin} was developed from
a theoretical analysis of the representational limits of GNNs within the
message-passing framework. Xu et al. proved that the expressive power of any
message-passing GNN is bounded by the Weisfeiler-Leman (WL) graph isomorphism
test, and demonstrated that GIN achieves this theoretical upper bound. The GIN
layer updates each atom by \textit{summing} its neighbours' representations, adding
the atom's own representation (scaled by a small learnable factor), and passing the
result through a multilayer perceptron.

The critical theoretical insight is that sum aggregation preserves multiset
information about neighbour features (unlike mean or max aggregation, which
can conflate distinct neighbourhood configurations), making GIN maximally
expressive among message-passing GNNs.

The GINEConv \parencite{hu2020strategies} extension incorporates edge features by
adding a transformed bond feature vector to each neighbour's representation before
the summation, so that bond-level chemistry informs the aggregation.

Despite its theoretical superiority, GIN does not always outperform architecturally
simpler GNNs on molecular property prediction benchmarks, suggesting that
the WL-level expressivity gains are primarily relevant for highly symmetric
graph structures that are uncommon in drug-like molecules. Empirically, GIN has
shown strong performance on graph classification tasks in the TU graph benchmarks
and MoleculeNet \parencite{wu2018moleculenet}.

\subsection{AttentiveFP}

AttentiveFP \parencite{xiong2020attentivefp} is a GNN purpose-built for molecular
property prediction, combining two mechanisms. At the \textbf{atom level}, attention
weights modulating each neighbour's contribution are computed from the two atoms'
representations \emph{together with the connecting bond's feature vector}, natively
integrating edge chemistry into the attention. At the \textbf{molecule level}, an
iterative readout refines a single molecular embedding: a context vector, initialised
from the mean atom representation, repeatedly attends over all atoms and is updated by
a Gated Recurrent Unit (GRU) over several timesteps, capturing progressively more
global structure. AttentiveFP reaches near state-of-the-art performance on several
MoleculeNet benchmarks \parencite{xiong2020attentivefp, jiang2021gnn}, which is why it
was chosen as the architecture for hyperparameter optimisation here---although the
benchmark in Chapter~4 ultimately found it the weakest of the five graph networks on
the scaffold test set.


\subsection{GraphSAGE}

GraphSAGE (Graph SAmple and aggreGatE) \parencite{hamilton2017graphsage} is an
inductive framework that learns aggregation functions from sampled local
neighbourhoods rather than the full adjacency matrix. Each atom forms the mean of its
neighbours' representations, concatenates this with its own previous representation,
and passes the result through a learnable transformation---preserving the atom's
identity alongside its local context. Mean aggregation is robust to high-degree atoms,
and the inductive formulation generalises naturally to novel scaffolds, which is
directly relevant under Bemis-Murcko evaluation. Here GraphSAGE uses BatchNorm and
concatenated mean+max pooling, following the same readout as GCN.

% ----------------------------------------------------------------
\section{Nepali Medicinal Plants and Natural Product Databases}
\label{sec:nepali_plants}
% ----------------------------------------------------------------

\subsection{Traditional Medicinal Plants in Nepal}

Nepal harbours remarkable biodiversity, hosting approximately 6,500 flowering plant
species within a geographic transect spanning tropical lowlands ($\sim$60 m a.s.l.)
to alpine zones ($>$5,000 m a.s.l.). Of these, over 700 species have documented
medicinal uses in traditional Ayurvedic and Tibetan medical systems.
Several Nepali medicinal plants have been
ethnopharmacologically associated with anticancer activity and are candidates
for systematic computational evaluation as potential EGFR inhibitor sources.

Seven species were selected on the basis of documented bioactivity, coverage in
COCONUT 2.0, and relevance to oncological conditions:
\begin{itemize}
    \item \textbf{\textit{Swertia chirayita}} (Chirayito) --- xanthones and secoiridoids;
    \item \textbf{\textit{Rhododendron arboreum}} (Lali Gurans) --- triterpenoids and flavonoids;
    \item \textbf{\textit{Nardostachys jatamansi}} (Jatamansi) --- sesquiterpenes;
    \item \textbf{\textit{Terminalia chebula}} (Haritaki) --- hydrolysable tannins, a Triphala constituent;
    \item \textbf{\textit{Berberis aristata}} (Chutro) --- the alkaloid berberine, with reported EGFR-pathway modulation;
    \item \textbf{\textit{Ocimum sanctum}} (Tulsi) --- eugenol and ursolic/rosmarinic acids, with reported EGFR-inhibitory activity;
    \item \textbf{\textit{Tinospora cordifolia}} (Giloy) --- diterpenoids and isoquinoline alkaloids.
\end{itemize}
All seven have ethnopharmacologically documented anti-inflammatory or antiproliferative use.

\subsection{The COCONUT 2.0 Database}

The \textbf{CO}llection of \textbf{O}pen \textbf{N}atural Prod\textbf{U}c\textbf{T}s
(COCONUT) is the largest freely available database of natural product structures,
curated and maintained by the Steinbeck group at Friedrich Schiller University Jena
\parencite{sorokina2021coconut}. COCONUT 2.0 contains over 400,000 unique natural
product structures drawn from more than 50 public natural product databases,
each associated with organism taxonomy, literature references, and precomputed
molecular descriptors. Structures are provided as canonical SMILES and InChI strings
with CAS registry numbers where available.

For this study, COCONUT 2.0 was used to retrieve SMILES for compounds annotated to
the seven Nepali plant species of interest using organism taxonomy matching. Retrieved
compounds were then subjected to the same preprocessing pipeline as the ChEMBL training
data (salt removal, SMILES canonicalisation, MW $\leq$ 1000 Da, RDKit validation) to
ensure featurisation consistency.

% ----------------------------------------------------------------
\section{Virtual Screening in Drug Discovery}
\label{sec:vs_lit}
% ----------------------------------------------------------------

\subsection{Overview of Virtual Screening}

Virtual screening (VS) is a computational strategy employed at the early stage
of drug discovery to prioritise a large library of candidate compounds for
experimental testing, based on predicted interaction with a biological target
or predicted biological activity. VS substantially reduces the number of compounds
that need to be synthesised and assayed experimentally, focusing experimental
resources on the most promising candidates \parencite{walters2020generative}.

Two principal paradigms of VS are distinguished:
\begin{itemize}
    \item \textbf{Structure-based virtual screening (SBVS)}: Requires a
    three-dimensional crystal or cryo-EM structure of the target protein. Molecular
    docking algorithms (AutoDock Vina, Glide, GOLD) predict the binding pose and
    estimate the binding affinity of each candidate compound within the target's
    binding site. SBVS is accurate but computationally intensive and requires
    high-resolution structural data.
    \item \textbf{Ligand-based virtual screening (LBVS)}: Does not require target
    structural knowledge. A model trained on known active and inactive compounds
    (e.g., from ChEMBL) is used to score and rank external compound libraries by
    predicted activity probability. LBVS is computationally scalable to millions
    of compounds and is particularly powerful when high-quality bioactivity training
    data is available.
\end{itemize}

This thesis implements a \textbf{ligand-based virtual screening} pipeline using its
two best-performing trained classifiers---the best graph neural network (GCN) and a
Random Forest baseline---as the scoring functions, and prioritising leads by
cross-model consensus rather than relying on any single model.

\subsection{GNN-Based Virtual Screening}

The landmark application of GNN-based virtual screening in drug discovery was
the discovery of \textbf{halicin} by Stokes et al. \parencite{stokes2020deep}
at the Collins Laboratory (MIT). A directed-message passing neural network (DMPNN)
was trained on 2,335 compounds from the Drug Repurposing Hub to predict broad-spectrum
antibiotic activity. The trained model was applied to score over 100 million compounds
from the ZINC15 database; halicin (SU-3327, a c-Jun N-terminal kinase inhibitor)
was identified as a top-ranked hit and subsequently validated as a potent,
broad-spectrum antibiotic in mouse models of \textit{Acinetobacter baumannii}
infection. This study demonstrated that GNN-based LBVS can identify structurally
novel bioactive compounds that escape detection by conventional similarity searches.

In the kinase inhibitor domain, GNN-based VS has been applied to several targets
including EGFR, CDK2, and VEGFR2, typically achieving hit rates of 10--50\% in
subsequent experimental validation---substantially higher than the 0.1--1\% hit
rates typical of unguided high-throughput screening \parencite{sun2021gcn}.

\subsection{Natural Product Libraries and Nepali Medicinal Plants}

A key methodological choice in virtual screening is the source and composition of
the screening library. This thesis employs a \textbf{natural product library} strategy:
bioactive compounds from seven traditional Nepali medicinal plant species are retrieved
from the COCONUT 2.0 database and used as the screening library. The biological and
ethnopharmacological rationale is well-established: plants such as
\textit{Swertia chirayita}, \textit{Berberis aristata}, and \textit{Ocimum sanctum}
have centuries of documented use in the treatment of cancer-related symptoms and have
yielded structurally diverse secondary metabolites---alkaloids, flavonoids, terpenes,
and tannins---that represent structurally distinct chemotypes compared to synthetic
kinase inhibitors, potentially offering novel scaffolds capable of engaging the EGFR
ATP-binding pocket through unconventional binding modes. This strategy directly
connects computational drug discovery to the biodiversity resources and traditional
knowledge of Nepal, supporting affordable and regionally accessible cancer therapeutic
research.

\subsection{ADMET Properties and Lead Prioritisation}

Virtual screening hit lists require secondary filtering by predicted ADMET
(Absorption, Distribution, Metabolism, Excretion, Toxicity) properties to
prioritise compounds with the greatest potential for clinical development
\parencite{walters2020generative}. In this study, ADMET-proxy properties were
computed using RDKit as physicochemical surrogates:

\begin{itemize}
    \item \textbf{Lipinski's Rule of Five (Ro5)}: MW $\leq$ 500 Da, LogP $\leq$ 5,
    HBD $\leq$ 5, HBA $\leq$ 10 as necessary (not sufficient) conditions for
    oral bioavailability \parencite{lipinski2001experimental}.
    \item \textbf{Quantitative Estimate of Drug-likeness (QED)}: A composite
    drug-likeness score (range 0--1) integrating MW, LogP, HBD, HBA, TPSA, number
    of aromatic rings, number of rotatable bonds, and the presence of structural
    alerts, calibrated to reflect the desirability profile of approved oral drugs
    \parencite{bickerton2012quantifying}.
    \item \textbf{Topological Polar Surface Area (TPSA)}: A molecular surface
    area descriptor correlated with oral absorption and blood-brain barrier
    permeability. Values $\leq$ 140 \AA$^2$ are associated with adequate
    gastrointestinal absorption.
\end{itemize}

% ----------------------------------------------------------------
\section{Research Gap}
\label{sec:research_gap}
% ----------------------------------------------------------------

A systematic analysis of the reviewed literature reveals three persistent and
interrelated methodological gaps in the application of Graph Neural Networks to
EGFR inhibitor discovery. These gaps collectively limit the clinical utility and
scientific rigour of existing computational approaches:

\begin{enumerate}
    \item \textbf{Evaluation Bias from Random Data Splitting.} The majority of
    published GNN models for EGFR or related kinase bioactivity prediction employ
    random train/test splitting \parencite{jiang2021gnn, malik2025deepegfr}.
    This practice allows molecules with near-identical Bemis-Murcko scaffolds to
    appear simultaneously in training and test partitions, creating an optimistic
    bias in reported performance metrics. A model that achieves 0.90 AUROC under
    random splitting may perform significantly worse when deployed on a genuinely
    novel scaffold---precisely the scenario encountered in prospective virtual
    screening. Scaffold-based splitting, which groups all molecules sharing the
    same core framework into a single partition, provides a more realistic and
    demanding assessment of model generalisation \parencite{wu2018moleculenet}.

    \item \textbf{Incomplete Molecular Featurisation.} Published GNN models for
    EGFR inhibitor classification frequently employ minimal node feature sets
    (e.g., atom type only, or a small number of atomic properties) and rarely
    incorporate explicit, multi-dimensional bond-level (edge) feature vectors
    \parencite{malik2025deepegfr}. Bond features encoding stereochemistry,
    conjugation, and aromaticity carry significant pharmacological information---for
    example, the aromatic bonds of the quinazoline hinge-binder scaffold and the
    acrylamide double bond of covalent inhibitors---that is discarded or only
    coarsely approximated by node-only models or ECFP fingerprints.

    \item \textbf{Lack of Prospective Screening on Natural Product and Regionally
    Relevant Compound Libraries.} Most academic GNN papers in the EGFR and kinase
    inhibitor domain report performance metrics on held-out test sets but do not
    deploy the trained model in a prospective virtual screening pipeline applied
    to natural product or ethnopharmacological compound collections
    \parencite{sun2021gcn, malik2025deepegfr}. For South Asian research communities
    where EGFR-mutant NSCLC is prevalent but novel targeted therapies are expensive
    and inaccessible, virtual screening of traditional medicinal plant compounds
    represents a high-impact translational application that is entirely absent from
    the computational EGFR inhibitor literature.
\end{enumerate}

\noindent Table~\ref{tab:gap_compare} formally positions this thesis against
representative published GNN studies for EGFR inhibitor prediction.

\begin{table}[H]
\centering
\caption{This thesis versus representative GNN-based EGFR studies across three methodological dimensions ($\checkmark$ addressed; $\circ$ partial; $\times$ absent).}
\label{tab:gap_compare}
{\singlespacing
\begin{tabular}{lccc}
\toprule
\multirow{2}{*}{\textbf{Study}} &
\textbf{Scaffold} & \textbf{Rich} & \textbf{Natural Product} \\
& \textbf{Split} & \textbf{Features} & \textbf{Screening} \\
\midrule
DeepEGFR \parencite{malik2025deepegfr}      & $\times$ & $\circ$ & $\times$ \\
GCN QSAR \parencite{jiang2021gnn}           & $\times$ & $\circ$ & $\times$ \\
AttentiveFP \parencite{xiong2020attentivefp} & $\times$ & $\checkmark$ & $\times$ \\
Hu et al. \parencite{hu2020strategies}      & $\circ$  & $\checkmark$ & $\times$ \\
\textbf{This Thesis}                        & $\checkmark$ & $\checkmark$ & $\checkmark$ \\
\bottomrule
\end{tabular}
}
\end{table}

This thesis addresses all three gaps through its two objectives, within a single
reproducible open-source pipeline. To close the evaluation and featurisation gaps
(1--2), it \textbf{benchmarks graph neural networks}---built with comprehensive
29-dimensional node and 12-dimensional edge features and tuned via Optuna---%
\textbf{against traditional machine-learning baselines} under a strict Bemis-Murcko
scaffold split. To close the screening gap (3), it \textbf{deploys the best-performing
models} in a prospective virtual screen of Nepali medicinal plant compounds from
COCONUT~2.0 with ADMET-proxy prioritisation.